\chapter{Conclusiones y trabajo futuro}
\label{cap:conclusiones}
Desarrollar proyectos con ideas más particulares y curiosas siempre ha sido algo que siempre me ha
gustado\footnote{Particulares y curiosas como el desarrollo de un simulador del procesador ficticio
\texttt{chip-8} para la Nintendo 3DS (\url{https://github.com/Luengor/Chip-8x3ds})}.
El concepto de detectar la postura del cuerpo humano y poder interactuar con un sistema informático
me parece algo muy interesante y que puede tener muchas aplicaciones interesantes, más allá de lo
útiles que puedan ser para la sociedad en general.

Trabajar en el proyecto ha sido una experiencia muy enriquecedora. He cementado y refinado mis
conocimientos sobre algunos aspectos que ya conocía como el desarrollo de \glspl{api}, el despliegue
de aplicaciones web con \texttt{Docker} o el uso del motor de videojuegos \texttt{Unity}. También
me ha permitido aprender nuevas tecnologías como todo lo relacionado con el \gls{frontend} de la
aplicación (\texttt{React}, \texttt{TypeScript}, \texttt{Bun}) y el uso de modelos de inteligencia
artificial a través de \texttt{TensorFlow.js}.

Fuera de las herramientas y tecnologías, he podido experimentar el proceso de desarrollo de un
proyecto desde cero, desde la concepción de la idea hasta la implementación final. En este aspecto,
los conocimientos adquiridos en la rama de la ingeniería del software han sido de gran ayuda,
especialmente en lo que respecta a la planificación y al diseño del proyecto. En cuanto al diseño
de la implementación, ha sido interesante ver los patrones de diseño y las arquitecturas de software
aplicados en el proyecto e implementados más allá de una simple descripción teórica.

Sin embargo, el proyecto no ha sido sin sus desafíos. Uno de los mayores retos ha sido integrar
todas las diferentes tecnologías y herramientas utilizadas en el proyecto. La comunicación entre
tres sistemas diferentes (\texttt{Unity}, \texttt{React} y \texttt{FastAPI}) en tres lenguajes
de programación distintos (\texttt{C\#}, \texttt{TypeScript} y \texttt{Python}) ha requerido un
esfuerzo considerable para garantizar que todo funcionara de manera correcta.

Otro desafío ha sido la implementación de la detección de posturas en tiempo real. Aunque una
versión simple se desarrolló con éxito en el primer sprint del proyecto, la implementación final
es el resultado de un proceso iterativo de refinamiento y mejora a lo largo de una parte muy
importante del proyecto.

Probablemente, el mayor desafío ha sido la gestión del tiempo y la planificación del proyecto.
Adaptarse a una nueva metodología de trabajo ha sido un proceso de prueba y error. Sprint tras
sprint, he ido aprendiendo a gestionar mejor el tiempo, a priorizar mejor las tareas y a calcular
correctamente el tiempo necesario para cada una de ellas. De cualquier manera, la planificación
ha dejado mucho que desear durante la totalidad del proyecto, y es algo que me gustaría mejorar
en futuros trabajos.

En cuanto al trabajo futuro, hay varias áreas en las que se podría mejorar el proyecto. Una de las
más importantes sería mejorar la precisión y la velocidad de la detección de posturas. Aunque
la implementación actual funciona bien, el área de la inteligencia artificial está en constante
evolución y siempre hay espacio para mejorar.

También se podría mejorar la interfaz de usuario y la experiencia del usuario. La interfaz actual
es algo pobre y podría beneficiarse de un diseño más atractivo. Funcionalidades como filtrar las
actividades realizadas por tipo, aumentar la personalización de la aplicación o añadir estadísticas
más completas sobre el rendimiento del usuario y su progreso son algunas de las mejoras que se
podrían implementar.

Por último, una zona de crecimiento bastante obvia es la posibilidad de añadir más juegos al limitado
repertorio de dos actividades que actualmente existen en la aplicación. El objetivo de este proyecto
se acercaba más a la creación de la plataforma y a la implementación de la detección de posturas
que a la creación de un amplio catálogo de juegos. La creación de nuevos juegos es un proceso
complejo, ya que requiere un diseño cuidadoso para garantizar que sean divertidos, pero su
implementación es relativamente sencilla una vez que se ha creado la infraestructura necesaria.

En resumen, el desarrollo de este trabajo de fin de grado ha sido una experiencia muy enriquecedora
y desafiante. He aprendido mucho sobre el desarrollo de software y he podido aplicar muchos de los
conocimientos adquiridos a lo largo de mi carrera. Aunque el proyecto no ha estado exento de
desafíos y problemas, he podido superarlos con creces (o por lo menos, he aprendido de ellos) y he
ganado una gran experiencia en el proceso.
